\documentclass[12pt]{article}

\usepackage{amsmath, amssymb, amsthm}
\usepackage{geometry}
\usepackage{hyperref}
\usepackage{enumitem}
\usepackage{setspace}
\usepackage{titlesec}

\geometry{margin=1in}
\setstretch{1.15}

\newtheorem{definition}{Definition}[section]
\newtheorem{proposition}{Proposition}[section]

\title{\textbf{The Kernel Architecture Framework}\\
\large Mathematical Core Specification}
\author{Gnome Badhi}
\date{May 24, 2026}

\begin{document}
\maketitle

\begin{abstract}
This document presents the mathematical core of the Kernel Architecture Framework.
It formalizes the kernel through its invariants, operator algebra, constraint manifold,
failure taxonomy, and application boundary. All content is expressed in minimal,
axiomatic form without explanatory prose.
\end{abstract}

\tableofcontents
\newpage

% ---------------------------------------------------------
% 1. PROVENANCE
% ---------------------------------------------------------
\section{Provenance}

\begin{definition}[Provenance Tuple]
\[
\mathcal{P}(K) = (a_0, t_0, \mathcal{C}_0, \mathcal{V})
\]
\end{definition}

\begin{proposition}[Provenance Preservation]
\[
\mathcal{P}(O(K)) = \mathcal{P}(K) \quad \forall O \in \mathcal{O}.
\]
\end{proposition}

% ---------------------------------------------------------
% 2. KERNEL ONTOLOGY
% ---------------------------------------------------------
\section{Kernel Ontology}

\subsection{Ontological Premise}

\begin{definition}[Kernel]
The kernel \(K\) is defined by its invariant set \(I(K)\) and admissible operator set \(\mathcal{O}\).
\end{definition}

\subsection{Primitive Elements}

\[
\sigma,\; \beta,\; O,\; s,\; C,\; \iota
\]

\subsection{Ontological Relations}

\[
O(s), \quad \beta(O), \quad P_i(s), \quad C(s)
\]

\subsection{Ontological Domains}

\[
D = \{D_{\mathrm{per}}, D_{\mathrm{con}}, D_{\mathrm{op}}, D_{\mathrm{ref}}\}
\]

\subsection{Identity Condition}

\[
\mathrm{Identity}(K) = I(K)
\]

\subsection{Boundary Condition}

\[
O(\beta) \subseteq \beta
\]

% ---------------------------------------------------------
% 3. KERNEL INVARIANTS
% ---------------------------------------------------------
\section{Kernel Invariants}

\subsection{Invariant Set}

\begin{definition}[Invariant]
\[
P_i : S \to \{0,1\}
\]
\end{definition}

\subsection{Invariant Classes}

\[
I(K) = I_{\mathrm{struct}} \cup I_{\mathrm{dyn}} \cup I_{\mathrm{sem}}
\]

\subsection{Preservation Rule}

\[
O \in \mathcal{O} \iff \forall P_i \in I(K),\; P_i(O(s)) = P_i(s)
\]

\subsection{Invariant Lattice}

\[
P_i \preceq P_j \iff P_j \Rightarrow P_i
\]

\subsection{Identity via Invariants}

\[
K_1 \equiv K_2 \iff I(K_1) = I(K_2)
\]

% ---------------------------------------------------------
% 4. OPERATOR ALGEBRA
% ---------------------------------------------------------
\section{Operator Algebra}

\subsection{Operator Space}

\[
\mathcal{O} = \{O : S \to S \mid O \text{ preserves } I(K)\}
\]

\subsection{Composition}

\[
O_2 \circ O_1 \in \mathcal{O}
\]

\subsection{Identity Operator}

\[
I_S(s) = s
\]

\subsection{Commutator}

\[
[O_i, O_j] = O_i O_j - O_j O_i
\]

\subsection{Operator Classes}

\[
\mathcal{O} =
\mathcal{O}_{\mathrm{struct}} \cup
\mathcal{O}_{\mathrm{dyn}} \cup
\mathcal{O}_{\mathrm{sem}}
\]

\subsection{Boundary-Respecting Condition}

\[
O(\beta) \subseteq \beta
\]

\subsection{Coherence Condition}

\[
C(O(s)) \ge C_{\min}
\]

% ---------------------------------------------------------
% 5. CONSTRAINT MANIFOLD
% ---------------------------------------------------------
\section{Constraint Manifold}

\subsection{Definition}

\[
\mathcal{M}_C = \bigcap_i \{s \in S \mid P_i(s) = 0\}
\]

\subsection{Forbidden Region}

\[
\mathcal{F} = S \setminus \mathcal{M}_C
\]

\subsection{Admissible Operator Condition}

\[
O(\mathcal{M}_C) \subseteq \mathcal{M}_C
\]

\subsection{Boundary of the Manifold}

\[
\partial \mathcal{M}_C = \{s \mid C(s) = C_{\min}\}
\]

\subsection{Stability}

\[
B_\epsilon(s) \subseteq \mathcal{M}_C
\]

\subsection{Constraint Gradient Field}

\[
G(s) = \sum_i \lambda_i \nabla P_i(s)
\]

\subsection{Constraint-Preserving Flow}

\[
\phi_t(s) \in \mathcal{M}_C \quad \forall t
\]

% ---------------------------------------------------------
% 6. FAILURE MODE TAXONOMY
% ---------------------------------------------------------
\section{Failure Mode Taxonomy}

\subsection{Primary Failures}

\[
\mathcal{F}_{\mathrm{primary}} =
\{\text{boundary},\; \text{invariant},\; \text{coherence},\; \text{operator},\; \text{cascade}\}
\]

\subsection{Secondary Failures}

\[
\mathcal{F}_{\mathrm{secondary}} =
\{\text{temporal},\; \text{semantic},\; \text{saturation},\; \text{overwrite}\}
\]

\subsection{Severity Functional}

\[
\mathrm{Severity}(s) =
E(s) + \alpha (C_{\min} - C(s))_+ + \beta \cdot \mathrm{CascadeDepth}(s)
\]

\subsection{Recovery Condition}

\[
\exists O \in \mathcal{O} : O(s) \in \mathcal{M}_C
\]

% ---------------------------------------------------------
% 7. APPLICATION BOUNDARY (MEANING-FUNCTIONAL VERSION)
% ---------------------------------------------------------
\section{Application Boundary}

Let $\mathcal{D}$ denote the space of all possible domains in which the kernel may be
applied. The application boundary partitions $\mathcal{D}$ into three operational classes:
valid, conditional, and forbidden. This partition is encoded by the meaning functional
\[
V : \mathcal{D} \to \{1, \tfrac12, 0\},
\]
where the codomain values are semantic response codes rather than numerical
quantities.

\subsection{Valid Region}

\[
\mathcal{A}_{\mathrm{valid}}
= \{ d \in \mathcal{D} \mid \mathcal{O}_d \subseteq \mathcal{O} \},
\qquad
V(d) = 1.
\]

\subsection{Conditional Region}

\[
\mathcal{A}_{\mathrm{conditional}}
= \{ d \in \mathcal{D} \mid \exists C_d : \mathcal{O}_d \cup C_d \subseteq \mathcal{O} \},
\qquad
V(d) = \tfrac12.
\]

\subsection{Forbidden Region}

\[
\mathcal{A}_{\mathrm{forbidden}}
= \mathcal{D} \setminus (\mathcal{A}_{\mathrm{valid}} \cup \mathcal{A}_{\mathrm{conditional}}),
\qquad
V(d) = 0.
\]

\subsection{Meaning-Functional Validity Test}

\[
V(d) =
\begin{cases}
1 & d \in \mathcal{A}_{\mathrm{valid}}, \\
\tfrac12 & d \in \mathcal{A}_{\mathrm{conditional}}, \\
0 & d \in \mathcal{A}_{\mathrm{forbidden}}.
\end{cases}
\]

\subsection{Boundary Conditions}

For any domain \(d\) with \(V(d) \neq 0\):
\[
\begin{aligned}
&\mathcal{O}_d \subseteq \mathcal{O}, \\
&C(O_d(s)) \ge C_{\min}, \\
&\mathcal{P}(O_d(K)) = \mathcal{P}(K).
\end{aligned}
\]
For \(V(d) = \tfrac12\), these conditions may require augmentation by domain-specific
constraints \(C_d\).

% ---------------------------------------------------------
% APPENDIX A — NOTATION
% ---------------------------------------------------------
\appendix
\section{Notation Reference}

\[
\begin{aligned}
&K &&\text{kernel} \\
&S &&\text{state space} \\
&s &&\text{state} \\
&I(K) &&\text{invariant set} \\
&\mathcal{O} &&\text{operator set} \\
&\mathcal{M}_C &&\text{constraint manifold} \\
&\mathcal{F} &&\text{forbidden region} \\
&C(s) &&\text{coherence functional} \\
&P_i &&\text{invariant predicate} \\
&G(s) &&\text{constraint gradient field} \\
&\phi_t &&\text{constraint-preserving flow}
\end{aligned}
\]

% ---------------------------------------------------------
% APPENDIX B — SYMBOL CLARIFICATION AND OCR ERROR MAP
% ---------------------------------------------------------
\section*{Appendix B — Symbol Clarification and OCR Error Map}

The following table lists common OCR or export substitutions observed in some
renderings of this manuscript. The left column shows the incorrect or ambiguous
form; the right column shows the correct mathematical symbol or expression.

\begin{center}
\begin{tabular}{ll}
\textbf{Incorrect / OCR Form} & \textbf{Correct Symbol or Expression} \\
\hline
\texttt{C} & $\subseteq$ \\
\texttt{U} & $\cup$ \\
\texttt{E} & $\in$ \\
\texttt{CO} & $\subseteq \mathcal{O}$ \\
\texttt{3} & $\exists$ \\
\texttt{V} & $\forall$ \\
\texttt{=>} & $\Rightarrow$ \\
\texttt{->} & $\to$ \\
\texttt{0 : S -> S} & $O : S \to S$ \\
\texttt{O2 001} & $O_2 \circ O_1$ \\
\texttt{aMc} & $\partial \mathcal{M}_C$ \\
\texttt{BE(s)} & $B_\epsilon(s)$ \\
\texttt{LA V P\_i(s)} & $\lambda_i \nabla P_i(s)$ \\
\texttt{Pt(s)} & $\phi_t(s)$ \\
\texttt{Avalid U Aconditional} & $\mathcal{A}_{\mathrm{valid}} \cup \mathcal{A}_{\mathrm{conditional}}$ \\
\texttt{V(d) = \{1,2,0\}} & $V(d) \in \{1, \tfrac12, 0\}$ \\
\texttt{Istruct U Idyn U Isem} & $I_{\mathrm{struct}} \cup I_{\mathrm{dyn}} \cup I_{\mathrm{sem}}$ \\
\end{tabular}
\end{center}

This table is authoritative for resolving any discrepancies between the intended
mathematical notation and artifacts introduced by PDF generation, OCR, or text
extraction systems.

\end{document}
